%------------------------------------------------------------------------------------
%  SECTION 7.1 : LE SYSTÈME EXCLUSIF D'ANNOTATIONS COLLABORATIVES TODOH
%------------------------------------------------------------------------------------
\section{Le Système Exclusif d'Annotations Collaboratives \texttt{\textbackslash todoH}}
\label{sec:systeme_todoh}

La phase de relecture entre étudiants, directeurs de thèse et pairs académiques est une étape cruciale.
Le paquet standard \textsf{todonotes} est souvent employé, mais il souffre de trois limitations majeures :
les messages longs débordent verticalement du bas de page, les caractères accentués UTF-8 génèrent des erreurs de compilation avec les paquets de soulignement traditionnels, et le texte surligné dans le corps n'offre aucun lien interactif vers le tableau de bord des tâches.

La classe \texttt{ClasseGuillaume} résout l'ensemble de ces verrous technologiques grâce à sa commande exclusive \texttt{\textbackslash todoH}.

\subsection{Les cinq innovations majeures de la macro \texttt{\textbackslash todoH}}
\label{subsec:cinq_innovations_todoh}

La syntaxe générale de la commande se décompose en quatre arguments :
\begin{center}
  \texttt{\textbackslash todoH\{options\_style\}\{couleur\}\{message\_marge\}\{texte\_surligné\}}
\end{center}

\begin{enumerate}
  \item \textbf{Surlignage robuste via \textsf{lua-ul} et \textsf{luacolor} :}
  Le texte cible est surligné dans la couleur du relecteur. Contrairement aux paquets obsolètes (\textsf{soul}), le moteur Lua gère les césures de mots, les accents français et les équations sans aucun bogue de compilation.
  
  \item \textbf{Troncature automatique des messages longs ($> 30$ caractères) :}
  Pour éviter qu'un commentaire volumineux ne vienne chevaucher la note suivante ou déborder de la marge, la classe mesure la longueur de la chaîne avec \textsf{xstring} et abrège la note à 30 caractères suivis de points de suspension.

  \item \textbf{Infobulle interactive (\textsf{pdfcomment}) :}
  Grâce à la commande \texttt{\textbackslash pdftooltip}, l'intégralité du message reste disponible : dans votre lecteur PDF (Adobe Acrobat, PDF-XChange, navigateurs modernes), \textbf{un simple survol de la souris sur la note dans la marge fait apparaître une bulle contextuelle} dévoilant le texte complet !

  \item \textbf{Détection intelligente des commandes et URL (\textsf{xurl}) :}
  Si le commentaire contient une balise \LaTeX{} (comme un lien hypertexte \texttt{\textbackslash url\{...\}}), la classe détecte la présence de la barre oblique inversée (\texttt{\textbackslash}) et suspend la troncature pour préserver l'intégrité du lien cliquable.

  \item \textbf{Lien hypertexte de retour vers la table des tâches :}
  Un clic direct sur le texte surligné dans le paragraphe téléporte instantanément le lecteur vers la liste récapitulative des révisions (\texttt{\textbackslash listoftodos}) située au début du mémoire.
\end{enumerate}

\subsection{Démonstrations concrètes en situation de révision}
\label{subsec:demos_todoh}

Voici deux exemples concrets illustrant le comportement adaptatif de la commande :

Dans ce premier passage, nous simulons une remarque de fond \todoH{FloraStyle}{red!30}{Cette démonstration théorique requiert une justification expérimentale plus détaillée avec des données chiffrées issues du banc d'essai.}{un raisonnement nécessitant un approfondissement méthodologique}.
Observez dans la marge de droite : la note est abrégée pour ne pas saturer l'espace, mais un survol à la souris révèle l'ensemble de la directive de relecture.

Dans ce second cas, nous transmettons à l'auteur une ressource bibliographique en ligne sans risque d'altération du lien : \todoH{GuillaumeStyle}{NavyBlue!30}{Vérifier les spécifications normatives sur \url{https://www.iso.org/standard/30669.html}}{la conformité avec les recommandations internationales ISO}.
Le lien URL est protégé et césuré proprement grâce au paquet \textsf{xurl}.

